\chapter{Implementing an Adapter}\label{ch:adapter}

An adapter makes your device an EtherNet/IP I/O device. A PLC opens a Class~1
connection to it and exchanges assembly data at a fixed interval, the RPI.
This chapter tells you how to start the adapter, exchange data, and set the
assemblies so that a PLC accepts them.

\section{Assemblies}\label{sec:assemblies}

An \term{assembly} is a block of bytes with an instance number. The adapter
has three types of assembly. Table~\ref{tab:asmtypes} gives their directions.

\begin{table}[htbp]
\centering
\caption{Assembly types}\label{tab:asmtypes}
\begin{tabularx}{\linewidth}{@{}P{0.25\linewidth}P{0.14\linewidth}L@{}}
\toprule
\textbf{Type} & \textbf{Direction} & \textbf{Meaning} \\
\midrule
\code{EIP_ASM_CONSUMED} & \OT & Data that the PLC writes. Your application reads it. \\
\code{EIP_ASM_PRODUCED} & \TO & Data that your application writes. The PLC reads it. \\
\code{EIP_ASM_CONFIG} & In the Forward\_Open & Configuration data. Size 0 is permitted. \\
\bottomrule
\end{tabularx}
\end{table}

\begin{caution}
Do not mix up the two directions. Consumed is what the PLC writes, and
produced is what the PLC reads. The built-in test set uses instance 100 as
consumed. The examples in this manual use instance 100 as produced.
\end{caution}

You select the instance numbers. This manual uses 150 (consumed, 32 bytes),
100 (produced, 32 bytes), and 151 (config, 0 bytes). If you give no
configuration, the adapter uses a built-in test set (table~\ref{tab:testset}).

\begin{table}[htbp]
\centering
\caption{The built-in test set (\code{config} or \code{assembly_count} is 0)}\label{tab:testset}
\begin{tabularx}{\linewidth}{@{}P{0.14\linewidth}P{0.20\linewidth}P{0.12\linewidth}L@{}}
\toprule
\textbf{Instance} & \textbf{Direction} & \textbf{Bytes} & \textbf{Notes} \\
\midrule
100 & Consumed & 8 & \\
101 & Produced & 8 & Initial data \code{11 22 33 44 55 66 77 88}. \\
102 & Produced & 4 & \\
150 & Config & 4 & \\
\bottomrule
\end{tabularx}
\end{table}

With \code{vendor_id} 0, the adapter also uses a test identity: vendor 65000
and product name ``EIP Test Adapter.''

\section{Start the Adapter}\label{sec:start}

\begin{procedure}{Start the built-in adapter}
\begin{enumerate}
  \item Declare an array of \code{EipAssemblyConfig}, one entry for each
        assembly.
  \item Declare an \code{EipAdapterConfig} and set all of it to zero with
        \code{memset}.
  \item Set the identity fields: \code{vendor_id}, \code{device_type},
        \code{product_code}, \code{revision_major}, \code{revision_minor}, and
        \code{product_name}.
  \item Set \code{assemblies} and \code{assembly_count}.
  \item Set the optional fields and callbacks (section~\ref{sec:callbacks}).
  \item Call \code{eip_adapter_start}.
  \item Make sure that the return value is \code{EIP_OK}.
  \item On a poll-mode port only, call \code{eip_adapter_poll} on each pass
        of the main loop.
\end{enumerate}
\end{procedure}

\begin{lstlisting}
#include <string.h>
#include "eip.h"

static const EipAssemblyConfig kAssemblies[] = {
    /* instance, direction,   size, name,     initial, implicit_only */
    {150, EIP_ASM_CONSUMED, 32, "Outputs", 0, 0},   /* O->T: the PLC writes it  */
    {100, EIP_ASM_PRODUCED, 32, "Inputs",  0, 0},   /* T->O: we send it         */
    {151, EIP_ASM_CONFIG,    0, "Config",  0, 0},   /* config point of the path */
};

int start_eip(void)
{
    EipAdapterConfig cfg;

    memset(&cfg, 0, sizeof cfg);
    cfg.vendor_id = 65000;            /* placeholder: use your ODVA vendor id */
    cfg.device_type = 43;             /* 43 generic device, 12 comms adapter */
    cfg.product_code = 1;
    cfg.revision_major = 1;
    cfg.revision_minor = 1;           /* an EDS needs 1 or higher */
    cfg.product_name = "My Device";   /* up to 31 characters */
    cfg.serial_number = 0;            /* 0: derived from the MAC */
    cfg.assemblies = kAssemblies;
    cfg.assembly_count = 3;
    return eip_adapter_start(&cfg);   /* copies what it needs from cfg */
}
\end{lstlisting}

\begin{caution}
Set all fields of \code{EipAdapterConfig}, including \code{link_info} and the
other callbacks, before you call \code{eip_adapter_start}. The function copies
the configuration. A change after the call has no effect.
\end{caution}

\begin{note}
\notelabel Vendor ID 65000 is a placeholder. Get your own vendor ID from
ODVA, and use it in the device and in its EDS file.
\end{note}

\code{eip_adapter_start} copies what it needs, so \code{cfg} and the assembly
array can be on the stack. The adapter listens on TCP 44818, UDP 44818, and
UDP 2222. If the network is not ready, the adapter tries again to open its
sockets one time each second. The field \code{listening} in
\code{EipAdapterView} becomes 1 when the sockets are open.

\section{Exchange Data}\label{sec:exchange}

Use \code{eip_adapter_read} to get the data that the PLC sent. Use
\code{eip_adapter_write} to set the data that the PLC receives.

\begin{lstlisting}
uint8_t out[32], in[32];
uint16_t len = sizeof out;

eip_adapter_read(150, out, &len);   /* latest data from the PLC          */
/* ... use out, fill in ... */
eip_adapter_write(100, in, 32);     /* sent on the next T->O packet      */
\end{lstlisting}

These rules apply:

\begin{itemize}
  \item The two calls are safe from any task.
  \item \code{eip_adapter_write} needs the exact size of the assembly.
  \item For \code{eip_adapter_read}, \code{data_len} is the capacity on input
        and the byte count on output. A capacity smaller than the assembly is
        \code{EIP_ERR_TOO_BIG}.
  \item An unknown instance returns \code{EIP_ERR_PATH}.
\end{itemize}

\subsection{The Safe State of Consumed Data}\label{sec:safestate}

The consumed assembly reads as all zeros in these conditions:

\begin{itemize}
  \item The connection timed out or closed.
  \item The PLC is in Program mode. The run/idle header then says ``idle.''
\end{itemize}

Before the first connection, the consumed assembly reads as its
\code{initial} data. That data is also zeros if you do not set it.
Figure~\ref{fig:safestate} shows the states.

\begin{figure}[htbp]
\centering
\begin{tikzpicture}[
  font=\sffamily\small,>=Stealth,
  st/.style={draw=steel,fill=white,rounded corners=6pt,minimum width=3.3cm,text width=3.1cm,minimum height=1.25cm,align=center,line width=0.8pt},
  zero/.style={st,fill=signal!6,draw=signal!70},
  live/.style={st,fill=teal!12,draw=teal},
  tr/.style={->,steel,line width=0.7pt},
  lab/.style={font=\sffamily\scriptsize,inner sep=2pt,align=center,text=navy}]
  \node[st] (init) at (0,0) {\textbf{No connection yet}\\\scriptsize reads \code{initial} data};
  \node[live] (run) at (5.3,0) {\textbf{Connected, run}\\\scriptsize reads live PLC data};
  \node[zero] (idle) at (11.6,0) {\textbf{Connected, idle}\\\scriptsize reads all zeros};
  \node[zero] (closed) at (5.3,-3.6) {\textbf{Timed out or closed}\\\scriptsize reads all zeros};
  \draw[tr] (init) -- node[lab,above]{Forward\_Open\\accepted} (run);
  \draw[tr] ([yshift=3mm]run.east) -- node[lab,above]{header says idle} ([yshift=3mm]idle.west);
  \draw[tr] ([yshift=-3mm]idle.west) -- node[lab,below]{header says run} ([yshift=-3mm]run.east);
  \draw[tr] ([xshift=-7mm]run.south) -- node[lab,left]{timeout or\\Forward\_Close} ([xshift=-7mm]closed.north);
  \draw[tr] ([xshift=7mm]closed.north) -- node[lab,right]{new\\Forward\_Open} ([xshift=7mm]run.south);
  \draw[tr] (idle.south) |- node[lab,pos=0.25,right]{timeout or\\Forward\_Close} (closed.east);
\end{tikzpicture}
\caption{States of the consumed assembly. Red states read as all zeros. A
modeless connection, which has no run/idle header, stays in ``run.''}
\label{fig:safestate}
\end{figure}

\begin{warning}
Design your application so that all-zero consumed data puts your machine
in a safe condition. When the connection stops, or when the PLC goes to
Program mode, your outputs see zeros. If zero is not a safe value for an
output, add logic that makes it safe. Do not use this function as a safety
function.
\end{warning}

\subsection{Status and Diagnostics}

\code{eip_adapter_snapshot} copies the adapter state into an
\code{EipAdapterView}. It gives the identity, the connection state, the RPI,
the inactivity timeout, and these counters:

\begin{itemize}
  \item packets in and out, and successful opens and timeouts
  \item drops by cause: old sequence number (\code{rx_dropped_seq}), wrong
        source address (\code{rx_dropped_source}), and malformed packet
        (\code{rx_dropped_bad})
  \item packets accepted in idle (\code{rx_idle}), and UDP replies not sent
        (\code{udp_replies_dropped})
\end{itemize}

Use these values on a status page. Appendix~\ref{app:types} lists all
fields. The Identity object status word also shows the state:
\hex{0034} with no connection, \hex{0065} in run, and \hex{0075} in idle.

\section{Assemblies, Sizes, and the EDS File}\label{sec:sizes}

The PLC sets up the connection from your EDS file. The adapter refuses a
Forward\_Open that does not agree with its assemblies.

\subsection{Assembly Size and Wire Size}

In an EDS file and in PLC software, you enter the \emph{assembly} size. The
2-byte sequence count and the 4-byte run/idle header are added on the wire.
The adapter checks the wire size. Let $n$ be the assembly size in bytes, and
let $h$ be 1 if the direction has a run/idle header, else 0. Then:
\begin{align}
  S_{\OT} &= 2 + 4h + n, \label{eq:o2t}\\
  S_{\TO} &= 2 + n. \label{eq:t2o}
\end{align}
For a 32-byte assembly in each direction, $S_{\OT} = 38$ and $S_{\TO} = 34$.
Figure~\ref{fig:o2t} shows where these bytes are in an \OT{} packet.

\begin{figure}[htbp]
\centering
\begin{tikzpicture}[font=\sffamily\scriptsize,
  f/.style={draw=steel,minimum height=1.0cm,align=center,inner sep=1pt,fill=white},
  g/.style={f,fill=steel!10},
  c/.style={f,fill=teal!14,draw=teal},
  d/.style={f,fill=teal!30,draw=teal}]
  \def\u{0.36}
  \node[g,minimum width=2*\u cm] (a) at (0,0) {Item\\count\\2};
  \node[g,minimum width=2*\u cm,right=0pt of a] (b) {Type\\\texttt{8002}\\2};
  \node[g,minimum width=2*\u cm,right=0pt of b] (c) {Len\\8\\2};
  \node[g,minimum width=4*\u cm,right=0pt of c] (e) {Connection\\ID\\4};
  \node[g,minimum width=4*\u cm,right=0pt of e] (h) {Sequence\\number\\4};
  \node[g,minimum width=2*\u cm,right=0pt of h] (i) {Type\\\texttt{00B1}\\2};
  \node[g,minimum width=2*\u cm,right=0pt of i] (j) {Len\\2};
  \node[c,minimum width=2.6*\u cm,right=0pt of j] (k) {Seq.\\count\\2};
  \node[c,minimum width=4*\u cm,right=0pt of k] (l) {Run/idle\\header\\4};
  \node[d,minimum width=7*\u cm,right=0pt of l] (m) {Assembly data\\$n$};
  \draw[decorate,decoration={brace,amplitude=4pt},steel] ([yshift=2pt]b.north west) -- node[above=4pt,align=center]{Sequenced Address item} ([yshift=2pt]h.north east);
  \draw[decorate,decoration={brace,amplitude=4pt},steel] ([yshift=2pt]i.north west) -- node[above=4pt,align=center]{Connected Data item} ([yshift=2pt]m.north east);
  \draw[decorate,decoration={brace,amplitude=4pt,mirror},teal] ([yshift=-2pt]k.south west) -- node[below=5pt,align=center]{Connection size in the Forward\_Open: $S_{\OT} = 2 + 4 + n$} ([yshift=-2pt]m.south east);
  \node[anchor=north west,text=steel] at ([yshift=-2pt]a.south west) {Common Packet Format, UDP 2222};
\end{tikzpicture}
\caption{An \OT{} Class~1 packet with a run/idle header. Sizes are in bytes.
A \TO{} packet has the same items but normally no run/idle header.}
\label{fig:o2t}
\end{figure}

The run/idle header is a format flag in the EDS file. It is not part of the
size that you enter.

\subsection{Rules for a Forward\_Open}\label{sec:forwardopen}

Table~\ref{tab:forwardopen} gives the rules that the adapter applies to a
Forward\_Open, and the status that it returns when a rule is not met. The
sizes in this table are wire sizes.

\begin{table}[htbp]
\centering
\caption{Forward\_Open rules and the error status for each}\label{tab:forwardopen}
\small
\begin{tabularx}{\linewidth}{@{}LP{0.30\linewidth}@{}}
\toprule
\textbf{Rule} & \textbf{Status if not met} \\
\midrule
The connection path names an existing consumed, produced, and config instance. & \hex{01} / \hex{012A} (consumed), \hex{012B} (produced), \hex{0129} (config) \\
\OT{} size $= 2 + 4h + n$ (equation~\ref{eq:o2t}). & \hex{01} / \hex{0127} \\
\TO{} size $= 2 + n$. If the size also includes a 4-byte header, the adapter accepts it and sends a header that always says ``run.'' & \hex{01} / \hex{0128} \\
Config data in the path, if any, is equal to the config assembly size. & \hex{01} / \hex{0126} \\
The electronic key, if the PLC checks it, agrees with the vendor, device type, product code, and revision. A zero in the key agrees with all values. & \hex{01} / \hex{0114}--\hex{0116} \\
No other originator owns the connection. & \hex{01} / \hex{0106}, or \hex{01} / \hex{0113} if the request has no \OT \\
\bottomrule
\end{tabularx}
\end{table}

A size-0 \code{EIP_ASM_CONFIG} assembly accepts a path with a config point
and no data. Most PLCs send that. Set \code{implicit_only} to 1 to hide an
assembly from explicit Get and Set (status \hex{16}). Class~1 can still use
that assembly.

\subsection{The Inactivity Timeout}\label{sec:timeout}

The adapter closes the connection if no valid \OT{} packet arrives in the
inactivity timeout. The timeout comes from the \OT{} RPI and the timeout
multiplier $m$ (0--7) in the Forward\_Open.\footnote{ODVA, \textit{The CIP
Networks Library}, vol.~1, \textit{Common Industrial Protocol}, chap.~3,
Connection Manager object.} Let $T_{\text{f}}$ be
\code{timeout_floor_ms} and let $\Delta$ be the clock resolution of the port:
\begin{equation}
  T_{\text{inactivity}} = \max\!\left(\text{RPI}_{\OT} \cdot 4 \cdot 2^{m},\;
  T_{\text{f}},\; 2\Delta\right).
  \label{eq:timeout}
\end{equation}
The result is rounded up to whole ticks. In ticks, it is also limited to
\hex{00FFFFFF}. For example, a Micro850 sends an RPI of 50~ms with $m = 0$,
so $T_{\text{inactivity}} = 50 \times 4 = 200$~ms. A multiplier of 2 with a
10~ms RPI gives $10 \times 16 = 160$~ms.

\begin{note}
\notelabel On a platform that is not real time, such as Windows, set
\code{timeout_floor_ms} to 500 or more. A 100~ms stall of the host then does
not close the connection.
\end{note}

\subsection{Multicast \TO}

If the PLC requests a multicast \TO{} connection, the adapter selects the
group address from its own IP address $a$ and its network mask $k$, as the
EtherNet/IP specification tells.\footnote{ODVA, \textit{The CIP Networks
Library}, vol.~2, \textit{EtherNet/IP Adaptation of CIP}, sec.~3-5.3.} Let the
host ID be $h = a \wedge \lnot k$. Then:
\begin{equation}
  A_{\text{mcast}} = 239.192.1.0 + 32 \cdot \bigl((h - 1) \wedge \hex{3FF}\bigr).
\end{equation}
Multicast does not let a second originator listen. The adapter still accepts
one Class~1 connection only.

\section{Optional Callbacks}\label{sec:callbacks}

Table~\ref{tab:callbacks} lists the callbacks in \code{EipAdapterConfig}. All
of them run in the adapter context, the task or \code{eip_adapter_poll}. Each
receives \code{callback_ctx} as its first argument.

\begin{table}[htbp]
\centering
\caption{Adapter callbacks}\label{tab:callbacks}
\small
\begin{tabularx}{\linewidth}{@{}P{0.17\linewidth}P{0.27\linewidth}L@{}}
\toprule
\textbf{Field} & \textbf{When it is called} & \textbf{Use} \\
\midrule
\code{link_info} & Each 100 ms & Gives link state, speed, and duplex for the Ethernet Link object, if the port cannot read them (ESP-IDF). Return 1 while the link is up, and 0 to report it down. \\
\code{apply_interface} & Approximately 150 ms after a scanner sets TCP/IP Interface attribute 3, 5, or 6. The reply has already gone out. & Store and apply the new IP settings. Set it to 0 to make those attributes read-only. \\
\code{reset} & Approximately 300 ms after an Identity Reset request & Restart (type 0), or set the factory settings and restart (type 1). Set it to 0 to refuse the service. \\
\code{allow_change} & Before the adapter accepts a TCP/IP Interface Set or an Identity Reset & Return 0 to permit the change. Return a CIP general status, usually \hex{10} (device state conflict), to refuse it. \\
\bottomrule
\end{tabularx}
\end{table}

\begin{warning}
A scanner on the network can send an Identity Reset or change the IP address
of your device. If your device controls machinery, supply
\code{allow_change}. Make it refuse these requests while the machinery can
move.
\end{warning}

\begin{caution}
\code{allow_change} runs while the adapter holds \code{eip_hal_lock}. Keep it
short. Do not call an \code{eip_adapter_*} function from it, because that
function also takes the lock.
\end{caution}

A request that \code{allow_change} refuses never gets to
\code{apply_interface} or \code{reset}.

\section{More Control: The EipServer API}\label{sec:eipserver}

The functions \code{eip_adapter_*} use one built-in instance of
\code{EipServer}. The header \file{eip_server.h} gives access to that type.
Use it for these tasks:

\begin{itemize}
  \item Put the 11~KB server in memory that you select.
  \item Run the loop from your own task.
  \item Stop the adapter and start it again.
\end{itemize}

\begin{lstlisting}
static EipServer server;

eip_server_init(&server, &config);
for (;;) {
    eip_server_service(&server, 5);   /* RTOS task: waits up to 5 ms */
}
/* bare metal: call eip_server_poll(&server) from the main loop instead */
\end{lstlisting}

\code{eip_server_init}, \code{eip_server_poll}, \code{eip_server_service},
and \code{eip_server_close} run in one context. \code{eip_server_read},
\code{eip_server_write}, and \code{eip_server_snapshot} are safe from any
task, because they take \code{eip_hal_lock}.

\subsection{The Sans-I/O Core}

The header \file{eip_adapter.h} gives the core under the server. The core
never uses a socket or a clock. You give it packets and a tick count, and you
send the packets that it returns. Use it only if you write your own event
loop. The core has one shared reply buffer. Thus you must serialize the calls
on all \code{EipAdapter} objects in a program. Appendix~\ref{app:api} lists
its functions.
